\section*{Chapitre \ref{ch:g12:stereochemistry} --- Stéréochimie~: chiralité et isomères}

\begin{solution}{exo:g12:stereochemistry:1}
Butan-2-ol~: le carbone 2 (\ce{H}, \ce{OH}, \ce{CH3}, \ce{CH2CH3}).
Propan-2-ol~: aucun (le carbone 2 porte deux \ce{CH3} identiques).
2-Chlorobutane~: le carbone 2. Acide lactique~: le carbone 2 (\ce{H},
\ce{OH},
\ce{CH3}, \ce{COOH}).
\end{solution}

\begin{solution}{exo:g12:stereochemistry:2}
Le premier~: les deux \ce{Cl} (prioritaires sur \ce{H}) sont du même
côté~: \omterm{def:g12:stereochemistry:z-e}{Z}. Le second~: à gauche, \ce{CH3} est prioritaire, côté haut~; à
droite, \ce{CH2CH3} est prioritaire, côté bas~: \omterm{def:g12:stereochemistry:z-e}{E}.
\end{solution}

\begin{solution}{exo:g12:stereochemistry:3}
Chiraux~: le gant, la vis. Non chiraux~: la cuillère, le cube (chacun
possède un plan de symétrie). \ce{CHFClBr} est chirale (quatre \omterm{def:g7:atoms-and-molecules:atom}{atomes}
différents sur le carbone)~; \ce{CH2Cl2} ne l'est pas.
\end{solution}

\begin{solution}{exo:g12:stereochemistry:4}
Un \omterm{def:g6:pure-substances-and-mixtures:mixture}{mélange} de quantités égales de deux \omterm{def:g12:stereochemistry:enantiomer}{énantiomères}. La carvone
racémique contient les deux \omterm{def:g12:stereochemistry:enantiomer}{énantiomères}~: le nez perçoit donc les deux
odeurs à la fois.
\end{solution}

\begin{solution}{exo:g12:stereochemistry:5}
Non~: le carbone 1 porte deux \omterm{def:g7:atoms-and-molecules:atom}{atomes} identiques, \ce{H} et \ce{H}~; les
échanger donne la même \omterm{def:g7:atoms-and-molecules:molecule}{molécule}.
\end{solution}

\begin{solution}{exo:g12:stereochemistry:6}
\setchemfig{atom sep=2em}
\chemfig{C(-[2]CH_2CH_3)(-[4]H_3C)(<[:-30]OH)(<:[:-150]H)} \quad miroir
\quad \chemfig{C(-[2]CH_2CH_3)(-[0]CH_3)(<[:-150]HO)(<:[:-30]H)}
\end{solution}

\begin{solution}{exo:g12:stereochemistry:7}
2-Chlorobutane~: 2 (un \omterm{def:g12:stereochemistry:chiral}{carbone asymétrique}). 2,3-Dichlorobutane~: 3
(deux \omterm{def:g12:stereochemistry:chiral}{carbones asymétriques} et deux moitiés identiques~: une paire
d'\omterm{def:g12:stereochemistry:enantiomer}{énantiomères} et une forme méso). Pent-2-ène~: 2 (\omterm{def:g12:stereochemistry:z-e}{Z} et \omterm{def:g12:stereochemistry:z-e}{E}).
\end{solution}

\begin{solution}{exo:g12:stereochemistry:8}
\omterm{def:g12:stereochemistry:enantiomer}{Énantiomères}~; \omterm{def:g12:stereochemistry:diastereomer}{diastéréoisomères}~; \omterm{def:g12:stereochemistry:diastereomer}{diastéréoisomères}~; la même \omterm{def:g7:atoms-and-molecules:molecule}{molécule}
(l'acide méso-tartrique est superposable à son image dans un miroir).
\end{solution}

\begin{solution}{exo:g12:stereochemistry:9}
Dans la \omterm{def:g12:stereochemistry:conformation}{conformation} décalée où les groupes \ce{CH3} sont opposés, les
gros groupes sont aussi éloignés que possible~: c'est la plus stable.
Dans la \omterm{def:g12:stereochemistry:conformation}{conformation} éclipsée, les deux \ce{CH3} se font face et se
gênent.
\end{solution}

\begin{solution}{exo:g12:stereochemistry:10}
Ses deux moitiés, chacune \ce{CH(OH)COOH}, sont images l'une de l'autre
dans un miroir~: un plan situé entre les deux carbones centraux coupe la
\omterm{def:g7:atoms-and-molecules:molecule}{molécule} en deux moitiés symétriques. Une \omterm{def:g7:atoms-and-molecules:molecule}{molécule} qui possède un plan
de symétrie est sa propre image dans un miroir~: elle n'est pas chirale.
\end{solution}

\begin{solution}{exo:g12:stereochemistry:11}
\setchemfig{atom sep=1.7em}
(\omterm{def:g12:stereochemistry:z-e}{Z})-pent-2-ène \chemfig{C(-[:150]H_3C)(-[:210]H)=C(-[:30]CH_2-CH_3)(-[:-30]H)}
et (\omterm{def:g12:stereochemistry:z-e}{E})-pent-2-ène \chemfig{C(-[:150]H_3C)(-[:210]H)=C(-[:30]H)(-[:-30]CH_2-CH_3)}.
\end{solution}

\begin{solution}{exo:g12:stereochemistry:12}
Comme pour l'acide tartrique, avec \ce{Br} à la place de \ce{OH} et
\ce{CH3} à la place de \ce{COOH}~: les deux \ce{Br} en triangle plein,
les deux en triangle hachuré (une paire d'\omterm{def:g12:stereochemistry:enantiomer}{énantiomères}), et un plein et
un hachuré, qui possède un plan de symétrie~: la forme méso. Trois
\omterm{def:g12:stereochemistry:stereoisomer}{stéréoisomères}.
\end{solution}

\begin{solution}{exo:g12:stereochemistry:13}
La moitié. Le chimiste peut séparer les deux \omterm{def:g12:stereochemistry:enantiomer}{énantiomères} (comme
Pasteur, par d'autres moyens), ou ne fabriquer que l'\omterm{def:g12:stereochemistry:enantiomer}{énantiomère} utile,
par une réaction elle-même chirale (un catalyseur \omterm{def:g12:stereochemistry:chiral}{chiral}, une enzyme, un
\omterm{def:g7:chemical-reactions:reactant}{réactif} de départ \omterm{def:g12:stereochemistry:chiral}{chiral}).
\end{solution}

\begin{solution}{exo:g12:stereochemistry:14}
Trois~: (2E,4E), (2Z,4Z) et (2E,4Z), ce dernier étant la même \omterm{def:g7:atoms-and-molecules:molecule}{molécule}
que (2Z,4E) lue depuis l'autre extrémité.
\end{solution}

\begin{solution}{exo:g12:stereochemistry:15}
$2^3 = 8$ \omterm{def:g12:stereochemistry:stereoisomer}{stéréoisomères}, soit 4 paires d'\omterm{def:g12:stereochemistry:enantiomer}{énantiomères}. Chaque
\omterm{def:g12:stereochemistry:stereoisomer}{stéréoisomère} a un \omterm{def:g12:stereochemistry:enantiomer}{énantiomère} et $8 - 2 = 6$ \omterm{def:g12:stereochemistry:diastereomer}{diastéréoisomères}.
\end{solution}

\begin{solution}{pb:g12:stereochemistry:1}
\textbf{1.} \ce{C4H6O6}; $M = 4 \times 12.0 + 6 \times 1.0 + 6 \times 16.0
= \qty{150.0}{g/mol}$.

\textbf{2.} Deux groupes carboxyle (acide) et deux groupes hydroxyle
(\omterm{def:g11:functional-groups:alcohol}{alcool}).

\textbf{3.} Les deux carbones centraux, liés chacun à \ce{H}, \ce{OH},
\ce{COOH} et \ce{CH(OH)COOH}.

\textbf{4.} $2^2 = 4$.

\textbf{5.} Tous deux en triangle plein.

\textbf{6.} Les deux \ce{OH} en triangle hachuré~; elle n'est pas
superposable à L~: c'est l'acide D-tartrique, son \omterm{def:g12:stereochemistry:enantiomer}{énantiomère}.

\textbf{7.} Après rotation autour de la liaison centrale, de sorte que
les deux \ce{OH} soient orientés du même côté, la \omterm{def:g7:atoms-and-molecules:molecule}{molécule} possède entre
ses deux carbones centraux un plan qui fait correspondre une moitié à
l'autre.

\textbf{8.} Non~: c'est l'acide méso-tartrique.

\textbf{9.} Des quatre combinaisons (plein, plein), (hachuré, hachuré),
(plein, hachuré), (hachuré, plein), les deux dernières sont la même
\omterm{def:g7:atoms-and-molecules:molecule}{molécule}, la forme méso, parce que les deux moitiés de la \omterm{def:g7:atoms-and-molecules:molecule}{molécule} sont
identiques.

\textbf{10.} L et D~: \omterm{def:g12:stereochemistry:enantiomer}{énantiomères}. L et méso~: \omterm{def:g12:stereochemistry:diastereomer}{diastéréoisomères}.

\textbf{11.} À \qty{169}{\celsius} également~: des \omterm{def:g12:stereochemistry:enantiomer}{énantiomères} ont les
mêmes propriétés physiques.

\textbf{12.} Non~: c'est un \omterm{def:g12:stereochemistry:diastereomer}{diastéréoisomère} de L, un composé différent
qui a ses propres propriétés.

\textbf{13.} Les récepteurs du nez sont chiraux et distinguent les
images dans un miroir~; un thermomètre, qui n'est pas \omterm{def:g12:stereochemistry:chiral}{chiral}, ne le peut
pas.

\textbf{14.} Non, c'est un \omterm{def:g6:pure-substances-and-mixtures:mixture}{mélange} de deux composés~; dans son ensemble
il n'est pas optiquement actif, les deux \omterm{def:g12:stereochemistry:enantiomer}{énantiomères} se compensant.

\textbf{15.} Un seul \omterm{def:g12:stereochemistry:enantiomer}{énantiomère}~: chaque cristal était formé de L seul
ou de D seul.

\textbf{16.} \qty{5.0}{g} de chaque.

\textbf{17.} La forme méso est un composé différent, qui ne fait pas
partie du \omterm{def:g12:stereochemistry:enantiomer}{mélange racémique} de L et de D.

\textbf{18.} En dissolvant chaque tas~: les solutions font tourner la
lumière polarisée d'un même angle en sens opposés (et les cristaux sont
images l'un de l'autre dans un miroir).

\textbf{19.} Quatre~: ses deux \omterm{def:g12:stereochemistry:chiral}{carbones asymétriques} portent des groupes
différents (\ce{Cl} sur l'un, \ce{OH} sur l'autre), donc aucune
combinaison ne possède de plan de symétrie, et aucune paire de
combinaisons ne donne la même \omterm{def:g7:atoms-and-molecules:molecule}{molécule}.

\textbf{20.} Trois~: L, D et méso.
\end{solution}
